% --- Содержимое файла materials/5_4.tex ---

\subsection{Теорема Банаха--Штейнгауза}

\begin{idea}[Смысл теоремы]
	\textbf{Принцип равномерной ограниченности.}
	Представьте, что у вас есть набор пружин (операторов). Вы проверяете их по одной: подвешиваете груз весом 1 кг, и ни одна пружина не растягивается до бесконечности. Подвешиваете 2 кг — тоже все держат.
	
	Теорема утверждает: если пружины «держат» любой конкретный вес (поточечная ограниченность), то существует \textbf{общий предел жесткости} для всех пружин сразу (равномерная ограниченность). Не может быть такого, что пружины становятся всё слабее и слабее, даже если ни одна не ломается полностью.
	
	\textbf{Важно:} Это работает только если пространство $X$ «жирное» (полное). В «дырявых» пространствах теорема ломается.
\end{idea}

\begin{theorem}[Банаха--Штейнгауза]\label{Banach-Steingaus}
	Пусть $X$ --- банахово пространство (полное линейное нормированное), $Y$ --- линейное нормированное пространство.
	Пусть $\mathcal{A} \subseteq \mathcal{L}(X, Y)$ --- семейство линейных непрерывных операторов.
	
	Если это семейство ограничено в каждой точке:
	\[ \forall x \in X \quad \sup_{A \in \mathcal{A}} \|Ax\|_Y < \infty, \]
	то оно ограничено равномерно (по норме операторов):
	\[ \sup_{A \in \mathcal{A}} \|A\| < \infty. \]
\end{theorem}

\begin{idea}[Идея доказательства]
	Мы разбиваем пространство $X$ на множества $X_n$ («области», где операторы не растягивают векторы сильнее, чем в $n$ раз).
	Так как для любой точки $x$ операторы ограничены, объединение всех $X_n$ дает всё пространство $X$.
	
	Здесь вступает в игру **Теорема Бэра о категориях**: полное пространство нельзя составить из «тощих» (нигде не плотных) множеств. Значит, хотя бы одно множество $X_m$ должно быть «жирным» — содержать внутри себя целый шар.
	Используя линейность, мы «разносим» ограниченность с этого шара на всё пространство.
\end{idea}


\begin{proof}
	\textbf{1. Построение множеств Лебега.}
	Рассмотрим множества, на которых семейство операторов ограничено числом $n$:
	\[ X_n = \left\{x \in X \mid \sup_{A \in \mathcal{A}} \|Ax\|_Y \leqslant n \right\}. \]
	Условие $\sup \|Ax\| \leqslant n$ эквивалентно тому, что $\forall A \in \mathcal{A}: \|Ax\| \leqslant n$. Поэтому:
	\[ X_n = \bigcap_{A \in \mathcal{A}} \{x \in X \mid \|Ax\|_Y \leqslant n\}. \]
	
	\textbf{2. Замкнутость $X_n$.}
	Для фиксированного $A$ множество $\{x \mid \|Ax\| \leqslant n\}$ является прообразом замкнутого шара $\overline{B}_n(0) \subset Y$ при непрерывном отображении $A$. Следовательно, оно замкнуто.
	Множество $X_n$ есть пересечение замкнутых множеств, значит, оно тоже \textbf{замкнуто}.
	
	\textbf{3. Применение теоремы Бэра.}
	По условию теоремы (поточечная ограниченность), для каждого $x$ супремум конечен, значит, каждый $x$ попадает в какое-то $X_n$.
	\[ X = \bigcup_{n=1}^{\infty} X_n. \]
	Так как $X$ — полное пространство, по \textbf{теореме Бэра} оно не может быть объединением счетного числа нигде не плотных множеств. Значит, существует номер $m$, для которого $X_m$ имеет внутреннюю точку.
	То есть существуют $x_0 \in X$ и $\varepsilon > 0$, такие что шар лежит в множестве:
	\[ \overline{B}_\varepsilon(x_0) \subset X_m. \]
	Это означает, что для любого $z \in \overline{B}_\varepsilon(x_0)$ и любого $A \in \mathcal{A}$ выполняется $\|Az\| \leqslant m$.
	
	\textbf{4. Оценка нормы операторов.}
	Возьмем произвольный $u \in X$ с единичной нормой ($\|u\|=1$).
	Представим его через векторы из шара: $x_0$ и $x_0 + \varepsilon u$.
	\[ A(u) = A\left( \frac{(x_0 + \varepsilon u) - x_0}{\varepsilon} \right) = \frac{1}{\varepsilon} (A(x_0 + \varepsilon u) - A(x_0)). \]
	Оценим норму, используя неравенство треугольника:
	\[ \|Au\| \leqslant \frac{1}{\varepsilon} \left( \|A(x_0 + \varepsilon u)\| + \|A(x_0)\| \right). \]
	Так как $x_0 \in X_m$ и $(x_0 + \varepsilon u) \in \overline{B}_\varepsilon(x_0) \subset X_m$, то нормы их образов не превосходят $m$:
	\[ \|Au\| \leqslant \frac{1}{\varepsilon} (m + m) = \frac{2m}{\varepsilon}. \]
	
	Эта оценка не зависит ни от $u$, ни от выбора оператора $A$.
	Следовательно, $\sup_{A \in \mathcal{A}} \|A\| \leqslant \frac{2m}{\varepsilon} < \infty$.
\end{proof}

\begin{corollary}[О предельном операторе]
	Пусть $X$ — банахово пространство, $Y$ — линейное нормированное.
	Пусть $\{A_n\} \subset \mathcal{L}(X, Y)$ — последовательность непрерывных операторов, которая сходится поточечно:
	\[ \forall x \in X \quad A_n x \to A x. \]
	Тогда предельный оператор $A$ также является \textbf{линейным и непрерывным} ($A \in \mathcal{L}(X, Y)$).
\end{corollary}

\begin{proof}
	Линейность предела очевидна. Докажем непрерывность (ограниченность).
	Так как последовательность $\{A_n x\}$ сходится, она ограничена для каждого $x$.
	По теореме Банаха-Штейнгауза, нормы $\|A_n\|$ ограничены в совокупности константой $C$.
	Тогда для любого $x$:
	\[ \|Ax\| = \lim_{n \to \infty} \|A_n x\| \leqslant \limsup_{n \to \infty} (\|A_n\| \cdot \|x\|) \leqslant C \|x\|. \]
	Значит, $A$ ограничен.
\end{proof}

\begin{remark}
	Теоремы о сходимости квадратурных формул и интерполяционных процессов (Билет №16) являются прямыми следствиями теоремы Банаха--Штейнгауза.
\end{remark}